\begin{figure}[!ht]
\centering
\begin{tikzpicture}[every node/.style={font=\small,align=center},alt={The ordinary moving file reconstructs target six using full frame four, delta five and delta six. An equivalent file with an additional full frame five reconstructs the same target using only full five and delta six.}]
\node[draw,minimum width=20mm,minimum height=10mm] (a) at (0,0) {4\\Full};
\node[draw,minimum width=20mm,minimum height=10mm] (b) at (3,0) {5\\Delta};
\node[draw,minimum width=20mm,minimum height=10mm] (c) at (6,0) {6\\Delta};
\draw[-{Latex}] (a)--(b);\draw[-{Latex}] (b)--(c);
\node[above=2mm of b] {Three records for the moving file};
\node[draw,minimum width=20mm,minimum height=10mm] (d) at (3,-1.6) {5\\Full};
\node[draw,minimum width=20mm,minimum height=10mm] (e) at (6,-1.6) {6\\Delta};
\draw[-{Latex}] (d)--(e);
\node[align=left] at (.1,-1.6) {An additional\\full frame};
\end{tikzpicture}
\caption{Two valid encodings reach identical target pixels through different replay starts.}
\label{fig:revision-lossless-replay}
\end{figure}
